\documentclass{article}
\usepackage{graphicx}
\usepackage{booktabs}
\usepackage{amsmath}
\usepackage{hyperref}

\title{Human-AI Teaming for Preference Data Construction: Effects on Downstream Model Behavior and Bias}
\author{[Intern Name], Spurthi Setty \\
Anote AI}
\date{}

\begin{document}

\maketitle

\begin{abstract}
\noindent
Preference datasets used for RLHF fine-tuning are increasingly constructed with AI assistance, yet the effects of human-AI collaboration protocols on downstream model behavior and bias remain poorly understood. We conduct a controlled 3-condition study comparing (A) human-only preference labeling, (B) LLM-only preference generation, and (C) human-AI collaborative labeling, where humans review and accept/reject AI-generated preference pairs. We fine-tune a base language model on each dataset using RLHF and evaluate on downstream task performance, bias metrics (demographic parity, representation bias), and red-team adversarial robustness. Our results show that human-AI collaboration achieves X\% of human-only quality at Y\% of the cost, while LLM-only preference data introduces significant anchoring bias that degrades red-team robustness.
\end{abstract}

\section{Introduction}
% RLHF depends on preference data quality
% AI-assisted annotation is cost-effective but understudied
% Research question: how does collaboration protocol affect downstream model?

\section{Related Work}
\subsection{Reinforcement Learning from Human Feedback}
\subsection{Human-AI Collaborative Annotation}
\subsection{Bias in Preference Data}

\section{Experimental Design}
\subsection{Conditions}
% A: Human-only (crowdworkers)
% B: LLM-only (GPT-4o)
% C: Human-AI collaborative (human reviews AI proposals)

\subsection{Preference Task}
% Instruction-following quality judgment

\subsection{Bias Metrics}
% Demographic parity, representation bias, anchoring bias score

\subsection{Red-Teaming Protocol}

\section{Experimental Setup}
\subsection{Base Model}
\subsection{RLHF Fine-Tuning Protocol}
\subsection{Evaluation Datasets}

\section{Results}
% Table: Per-condition task performance + bias metrics
% Figure: Bias distribution across conditions
% Figure: Red-team success rates by condition

\subsection{Task Performance}
\subsection{Bias Analysis}
\subsection{Red-Teaming Results}

\section{Analysis}
\subsection{Anchoring Bias in LLM-only Data}
\subsection{Cost-Quality Trade-off}

\section{Conclusion}

\section*{References}
\small
\begin{enumerate}
    \item Christiano et al. Deep Reinforcement Learning from Human Preferences, 2017. NeurIPS 2017.
    \item Ouyang et al. Training Language Models to Follow Instructions with Human Feedback, 2022.
    \item Gilardi et al. ChatGPT Outperforms Crowd Workers for Text-Annotation Tasks, 2023. PNAS.
    % Add more references
\end{enumerate}

\end{document}
